% Paket & pengaturan global. Dimuat oleh main-standalone.tex / main-appendix.tex.

% Buang tanggal/ID yang berubah tiap build supaya PDF yang di-commit hanya
% berubah ketika isinya benar-benar berubah.
\pdfinfoomitdate=1
\pdftrailerid{}

\usepackage[T1]{fontenc}
\usepackage[utf8]{inputenc}
\usepackage{lmodern}
\usepackage{microtype}
\usepackage{textcomp}

\usepackage[english,indonesian]{babel}

\usepackage[a4paper,margin=2.5cm,headheight=15pt,footskip=1.2cm]{geometry}

\usepackage{amsmath}
\usepackage{graphicx}
\usepackage{adjustbox}
\usepackage{float}
\usepackage{caption}
\usepackage{booktabs}
\usepackage{longtable}
\usepackage{array}
\usepackage{multirow}
\usepackage{makecell}
\usepackage{enumitem}
\usepackage{xcolor}
\usepackage{tcolorbox}
\usepackage{listings}
\usepackage{fancyhdr}
\usepackage{titlesec}
\usepackage{needspace}
\usepackage{setspace}
\usepackage{tocloft}
\usepackage{pdflscape}
\usepackage{etoolbox}
\usepackage{xstring}

% --- TikZ & friends --------------------------------------------------------
\usepackage{tikz}
\usetikzlibrary{
  babel,            % wajib: babel indonesian mengaktifkan ; : ! sebagai shorthand
  arrows.meta, positioning, shapes.geometric, shapes.misc, shapes.arrows,
  fit, backgrounds, calc, decorations.pathreplacing, patterns, chains
}
\usepackage{circuitikz}
\usepackage{pgfplots}
\pgfplotsset{compat=1.18}

% --- Tampilan --------------------------------------------------------------
\captionsetup{
  font=small, labelfont={bf,color=accent}, labelsep=period,
  justification=raggedright, singlelinecheck=false, skip=6pt
}
\setlength{\parskip}{0.45em}
\setlength{\parindent}{0pt}
\onehalfspacing

\titleformat{\chapter}[hang]
  {\normalfont\LARGE\bfseries\color{navy}}
  {\thechapter.}{0.6em}{}
\titlespacing*{\chapter}{0pt}{-10pt}{18pt}
\titleformat{\section}{\normalfont\large\bfseries\color{navy}}{\thesection}{0.6em}{}
\titleformat{\subsection}{\normalfont\normalsize\bfseries\color{ink}}{\thesubsection}{0.6em}{}

\pagestyle{fancy}
\fancyhf{}
\fancyhead[L]{\footnotesize\color{muted}\nouppercase{\leftmark}}
\fancyhead[R]{\footnotesize\color{muted}\reporttitle}
\fancyfoot[C]{\footnotesize\color{muted}\thepage}
\renewcommand{\headrulewidth}{0.3pt}
\renewcommand{\headrule}{\color{accentMid}\hrule width\headwidth height\headrulewidth}

% --- Kode ------------------------------------------------------------------
\lstdefinestyle{report}{
  backgroundcolor=\color{codeBg},
  basicstyle=\ttfamily\footnotesize\color{ink},
  keywordstyle=\color{accent}\bfseries,
  commentstyle=\color{muted}\itshape,
  stringstyle=\color{warn},
  numbers=none, frame=single, rulecolor=\color{accentMid},
  framesep=5pt, xleftmargin=6pt, xrightmargin=6pt,
  breaklines=true, breakatwhitespace=false,
  postbreak=\mbox{\textcolor{muted}{$\hookrightarrow$}\space},
  columns=fullflexible, keepspaces=true, showstringspaces=false,
  aboveskip=8pt, belowskip=8pt
}
\lstset{style=report}

% --- Kotak catatan ---------------------------------------------------------
\newtcolorbox{notebox}[1][]{
  colback=accentLight, colframe=accent, boxrule=0.4pt, arc=2pt,
  left=8pt, right=8pt, top=6pt, bottom=6pt,
  fonttitle=\bfseries\small, title={Catatan}, #1
}
\newtcolorbox{warnbox}[1][]{
  colback=warnBg, colframe=warn, boxrule=0.4pt, arc=2pt,
  left=8pt, right=8pt, top=6pt, bottom=6pt,
  fonttitle=\bfseries\small, title={Peringatan}, #1
}

% --- Hyperref (paling akhir) -----------------------------------------------
\usepackage{hyperref}
\hypersetup{
  colorlinks=true,
  linkcolor=accent,
  citecolor=accent,
  urlcolor=accent,
  filecolor=accent,
  bookmarksnumbered=true,
  bookmarksopen=true
}
\usepackage{bookmark}
